\begin{proof}[Beweis des Satzes von Jordan für reell-analytische Kurven (Strahlenparitäts-Ansatz)\label{jordan:proof:jordan_reell_analytische_kurven_strahlenparitaet}]
    Für $p \in \mathbb{R}^2 \setminus \gamma$ wählen wir einen Strahl $L$ mit Anfangspunkt $p$, der nicht tangential an $\gamma$ verläuft und durch keinen degenerierten Punkt von $\gamma$ geht.
    Nach dem Lemma der endlichen Schnittstruktur ist die Anzahl der Schnittpunkte $N(p,L)$ diskret, also ist jede Richtung zulässig. 
    Sei N(p,L) die Anzahl der Schnittpunkte von $L \cap \gamma$. Nach dem Satz der endlichen Schnittstruktur ist ist diese Zahl endlich. 
    Die Parität $N(p,L) \mod 2$ ist wohldefiniert und hängt nicht von der Wahl des Strahls ab. 
    Der Wegzusammenhang von $I$ und $A$ und die Beschränktheit von $I$ ergeben sich aus der Konstruktion der tubularen Umgebung. 
    Die Beweise der Offenheit von $I$ und $A$, genauso wie die Unbeschränktheit von $A$, haben sich nicht verändert.
    Daher gilt der Satz von Jordan auch für reell-analytische, einfache geschlossene Kurven in der Ebene $\mathbb{R}^2$.
\end{proof}